% Options for packages loaded elsewhere
\PassOptionsToPackage{unicode}{hyperref}
\PassOptionsToPackage{hyphens}{url}
\PassOptionsToPackage{space}{xeCJK}
%
\documentclass[
  aps,
  onecolumn,
  11pt]{revtex4-2}
\usepackage{amsmath,amssymb}
\usepackage{iftex}
\ifPDFTeX
  \usepackage[T1]{fontenc}
  \usepackage[utf8]{inputenc}
  \usepackage{textcomp} % provide euro and other symbols
\else % if luatex or xetex
  \usepackage{unicode-math} % this also loads fontspec
  \defaultfontfeatures{Scale=MatchLowercase}
  \defaultfontfeatures[\rmfamily]{Ligatures=TeX,Scale=1}
\fi
\usepackage{lmodern}
\ifPDFTeX\else
  % xetex/luatex font selection
  \setmainfont[]{Times New Roman}
  \setmonofont[]{Menlo}
  \ifXeTeX
    \usepackage{xeCJK}
    \setCJKmainfont[]{Songti SC}
          \fi
  \ifLuaTeX
    \usepackage[]{luatexja-fontspec}
    \setmainjfont[]{Songti SC}
  \fi
\fi
% Use upquote if available, for straight quotes in verbatim environments
\IfFileExists{upquote.sty}{\usepackage{upquote}}{}
\IfFileExists{microtype.sty}{% use microtype if available
  \usepackage[]{microtype}
  \UseMicrotypeSet[protrusion]{basicmath} % disable protrusion for tt fonts
}{}
\makeatletter
\@ifundefined{KOMAClassName}{% if non-KOMA class
  \IfFileExists{parskip.sty}{%
    \usepackage{parskip}
  }{% else
    \setlength{\parindent}{0pt}
    \setlength{\parskip}{6pt plus 2pt minus 1pt}}
}{% if KOMA class
  \KOMAoptions{parskip=half}}
\makeatother
\usepackage{xcolor}
\usepackage{color}
\usepackage{fancyvrb}
\newcommand{\VerbBar}{|}
\newcommand{\VERB}{\Verb[commandchars=\\\{\}]}
\DefineVerbatimEnvironment{Highlighting}{Verbatim}{commandchars=\\\{\}}
% Add ',fontsize=\small' for more characters per line
\newenvironment{Shaded}{}{}
\newcommand{\AlertTok}[1]{\textcolor[rgb]{1.00,0.00,0.00}{\textbf{#1}}}
\newcommand{\AnnotationTok}[1]{\textcolor[rgb]{0.38,0.63,0.69}{\textbf{\textit{#1}}}}
\newcommand{\AttributeTok}[1]{\textcolor[rgb]{0.49,0.56,0.16}{#1}}
\newcommand{\BaseNTok}[1]{\textcolor[rgb]{0.25,0.63,0.44}{#1}}
\newcommand{\BuiltInTok}[1]{\textcolor[rgb]{0.00,0.50,0.00}{#1}}
\newcommand{\CharTok}[1]{\textcolor[rgb]{0.25,0.44,0.63}{#1}}
\newcommand{\CommentTok}[1]{\textcolor[rgb]{0.38,0.63,0.69}{\textit{#1}}}
\newcommand{\CommentVarTok}[1]{\textcolor[rgb]{0.38,0.63,0.69}{\textbf{\textit{#1}}}}
\newcommand{\ConstantTok}[1]{\textcolor[rgb]{0.53,0.00,0.00}{#1}}
\newcommand{\ControlFlowTok}[1]{\textcolor[rgb]{0.00,0.44,0.13}{\textbf{#1}}}
\newcommand{\DataTypeTok}[1]{\textcolor[rgb]{0.56,0.13,0.00}{#1}}
\newcommand{\DecValTok}[1]{\textcolor[rgb]{0.25,0.63,0.44}{#1}}
\newcommand{\DocumentationTok}[1]{\textcolor[rgb]{0.73,0.13,0.13}{\textit{#1}}}
\newcommand{\ErrorTok}[1]{\textcolor[rgb]{1.00,0.00,0.00}{\textbf{#1}}}
\newcommand{\ExtensionTok}[1]{#1}
\newcommand{\FloatTok}[1]{\textcolor[rgb]{0.25,0.63,0.44}{#1}}
\newcommand{\FunctionTok}[1]{\textcolor[rgb]{0.02,0.16,0.49}{#1}}
\newcommand{\ImportTok}[1]{\textcolor[rgb]{0.00,0.50,0.00}{\textbf{#1}}}
\newcommand{\InformationTok}[1]{\textcolor[rgb]{0.38,0.63,0.69}{\textbf{\textit{#1}}}}
\newcommand{\KeywordTok}[1]{\textcolor[rgb]{0.00,0.44,0.13}{\textbf{#1}}}
\newcommand{\NormalTok}[1]{#1}
\newcommand{\OperatorTok}[1]{\textcolor[rgb]{0.40,0.40,0.40}{#1}}
\newcommand{\OtherTok}[1]{\textcolor[rgb]{0.00,0.44,0.13}{#1}}
\newcommand{\PreprocessorTok}[1]{\textcolor[rgb]{0.74,0.48,0.00}{#1}}
\newcommand{\RegionMarkerTok}[1]{#1}
\newcommand{\SpecialCharTok}[1]{\textcolor[rgb]{0.25,0.44,0.63}{#1}}
\newcommand{\SpecialStringTok}[1]{\textcolor[rgb]{0.73,0.40,0.53}{#1}}
\newcommand{\StringTok}[1]{\textcolor[rgb]{0.25,0.44,0.63}{#1}}
\newcommand{\VariableTok}[1]{\textcolor[rgb]{0.10,0.09,0.49}{#1}}
\newcommand{\VerbatimStringTok}[1]{\textcolor[rgb]{0.25,0.44,0.63}{#1}}
\newcommand{\WarningTok}[1]{\textcolor[rgb]{0.38,0.63,0.69}{\textbf{\textit{#1}}}}
\setlength{\emergencystretch}{3em} % prevent overfull lines
\providecommand{\tightlist}{%
  \setlength{\itemsep}{0pt}\setlength{\parskip}{0pt}}
\setcounter{secnumdepth}{-\maxdimen} % remove section numbering
\ifLuaTeX
  \usepackage{selnolig}  % disable illegal ligatures
\fi
\IfFileExists{bookmark.sty}{\usepackage{bookmark}}{\usepackage{hyperref}}
\IfFileExists{xurl.sty}{\usepackage{xurl}}{} % add URL line breaks if available
\urlstyle{same}
\hypersetup{
  pdftitle={Cyclotron Frequency Comparison User Guide},
  pdfauthor={TrapLab documentation},
  hidelinks,
  pdfcreator={LaTeX via pandoc}}

\title{Cyclotron Frequency Comparison User Guide}
\author{TrapLab documentation}
\date{2026-10-24}

\begin{document}
\maketitle

Use the existing Lab cyclotron calculator to compare one to ten
explicitly identified ions or isotope-resolved molecules at a shared
magnetic field. Each row retains its own mass convention, uncertainty
and source. This tool calculates free cyclotron frequencies; it does not
fit a resonance, calibrate a magnetic field or infer an atomic
composition.

\subsection{Compare two ions on the
website}\label{compare-two-ions-on-the-website}

\begin{enumerate}
\def\labelenumi{\arabic{enumi}.}
\tightlist
\item
  Open Lab and find the cyclotron calculator. Keep \(B=7\) T for the
  first comparison. Leave field uncertainty blank to retain unknown
  uncertainty, or enter zero only when you deliberately treat the field
  as fixed.
\item
  Enter \texttt{133Cs} in the first row and select evaluated masses. Add
  a species row containing \texttt{85Rb}. Leave charge blank for the
  parsed charge or the default \(+1\), or enter a consistent nonzero
  integer explicitly.
\item
  Choose \textbf{Saved AME2020 / NUBASE2020} for a reproducible offline
  comparison, then \textbf{Calculate all rows}. Each row should show
  frequency in Hz, charge, ion mass and either a known uncertainty or an
  explicit unknown state. The Cs frequency is approximately 0.81 MHz at
  7 T; use the returned full precision rather than this rounded check
  value.
\item
  Open each row's source details. Confirm the evaluation, nuclear-state
  identifier, source origin and mass uncertainty. Online-first may fall
  back explicitly to the saved 2020 edition; online-only reports failure
  if its ground-state request cannot be fulfilled.
\item
  Save all rows JSON and download results CSV. Change a row and observe
  that its previous result becomes stale. Load the JSON again and verify
  recalculated results and retained source details. These files go to
  the browser's selected download location.
\end{enumerate}

Add permits at most ten rows; Remove leaves at least one. A field edit
invalidates affected results, and newer input cancels older asynchronous
fills. Do not copy a displayed number after changing its input without
recalculating.

\subsection{Species and manual mass
entry}\label{species-and-manual-mass-entry}

Accepted examples include \texttt{133Cs}, \texttt{Cs-133},
\texttt{cesium-133}, \texttt{133Xe{[}m1{]}}, \texttt{12C:16O*2} and
\texttt{{[}133Cs+1H{]}+}. Every molecular component must have an
isotope. Bare \texttt{Cs}, \texttt{H2O}, unavailable states and
conflicting charge indications are rejected. Use a bracketed explicit
charge or the charge field when a trailing digit could be interpreted as
an isotope count.

For a user mass, clear evaluated-mass selection and enter mass and
optional uncertainty. Choose \textbf{Neutral atom or component sum} or
\textbf{Already ionic mass} explicitly. Units are u or keV/\(c^2\). The
keV/\(c^2\) field is a full mass, not a mass excess. Changing its unit
converts both mass and known uncertainty. Editing an automatically
filled mass clears the former evaluated uncertainty and switches to
manual entry, preventing an old error from being attached to a new
value.

Charge is a nonzero integer between −100 and 100. The field and
resulting ion mass must be positive. The reported frequency is a
magnitude; signed frequency is retained in the exported data.

\subsection{Formula and units}\label{formula-and-units}

Let \(M\) and \(m_e\) be numerical masses in u and let \(u\) denote the
atomic mass constant in kg. The entered convention controls whether the
electron correction is applied:

\[m_{\mathrm{ion}}=M-qm_e\quad\text{for neutral input},\]

\[m_{\mathrm{ion}}=M\quad\text{for already ionic input}.\]

For positive ions the first expression subtracts electron masses; for
negative ions it adds them. Ionization, electron-affinity and molecular
binding energies are omitted. Already ionic input is not corrected
twice.

\[\nu_c=\frac{|q|eB}{2\pi m_{\mathrm{ion}}u},\qquad \nu_{\mathrm{signed}}=\frac{qeB}{2\pi m_{\mathrm{ion}}u}.\]

Frequency is in Hz for \(B\) in T. The mass-energy conversion retained
by this implementation is

\[C_u=uc^2=931494.10242\ {\mathrm{keV}},\]

\[M=\frac{E_{\mathrm{mass}}}{C_u},\qquad M=A+\frac{ME}{C_u}.\]

The first relation converts a full mass entered in keV/\(c^2\); the
second explains conversion of an evaluated atomic mass excess \(ME\).
They are different inputs and must not be interchanged.

\subsection{Uncertainty and
correlations}\label{uncertainty-and-correlations}

When both mass and field uncertainties are supplied, the displayed
marginal uncertainty is

\[\sigma_{\nu}=\nu_c\sqrt{\left(\frac{\sigma_B}{B}\right)^2+\left(\frac{\sigma_M}{m_{\mathrm{ion}}}\right)^2}.\]

Constants are held fixed. Blank uncertainty is unknown; numerical zero
is a known fixed input. If either required uncertainty is unknown,
frequency uncertainty remains unknown. For a repeated isotope in a
molecule, primitive coefficients are combined before propagation:

\[M_{\mathrm{sum}}=\sum_i n_iM_i,\qquad \sigma_{M_{\mathrm{sum}}}^2=\sum_i n_i^2\sigma_{M_i}^2.\]

This is a diagonal primitive model, not the complete AME covariance.
Rows share the same field. Its contribution to their covariance is

\[\operatorname{Cov}_B(\nu_j,\nu_k)=\frac{\nu_j\nu_k}{B^2}\sigma_B^2.\]

Shared isotope primitives can correlate rows as well. The calculator
supplies row marginals; it does not provide a joint frequency-ratio fit
or justify treating all displayed rows as independent.

\subsection{Evaluations and website versus native
conventions}\label{evaluations-and-website-versus-native-conventions}

Saved ground-state masses come from AME2020; exact state identities and
excitation information come from NUBASE2020. In the website convention,

\[M^*=M_{g,\mathrm{AME2020}}+\frac{E_{x,\mathrm{NUBASE2020}}}{C_u}.\]

The web online choice queries IAEA LiveChart ground-state mass excess
while retaining the explicit 2020 source policy and saved fallback.
Source evaluation date and request/retrieval date are distinct. The
native \texttt{ame2020} policy uses AME ground masses and direct NUBASE
selected-state masses, rather than reconstructing excited masses with
the above expression. Native \texttt{nubase2020} explicitly uses that
saved evaluation. This distinction can change an isomer result; preserve
the state and source policy with every exported comparison.

The web keeps its CODATA2018/AME2020 constant profile, including
\(m_e=5.48579909065\times10^{-4}\) u and
\(u=1.66053906660\times10^{-27}\) kg. The native compatibility profile
retains the same mass-energy conversion but uses
\(m_e=5.485799090441\times10^{-4}\) u and
\(u=1.66053906892\times10^{-27}\) kg. Small cross-platform numerical
differences therefore do not necessarily indicate a bug.

\subsection{Run the installed
companion}\label{run-the-installed-companion}

Open a terminal in the separately installed companion application
folder. For a saved-source command-line comparison, run

\begin{Shaded}
\begin{Highlighting}[]
\ExtensionTok{./Run\_private.command} \AttributeTok{{-}m}\NormalTok{ traplab\_qt.cyclotron\_compare 133Cs 85Rb }\AttributeTok{{-}{-}B}\NormalTok{ 7 }\AttributeTok{{-}{-}sigma{-}B}\NormalTok{ 0 }\AttributeTok{{-}{-}policy}\NormalTok{ ame2020}
\end{Highlighting}
\end{Shaded}

Append \texttt{-\/-output\ cyclotron-example} to create a new bundle
directory containing \texttt{cyclotron\_rows.json} and
\texttt{cyclotron\_results.csv}. An existing directory is rejected,
preserving earlier results. Use
\texttt{-\/-settings\ cyclotron-example/cyclotron\_rows.json} to reopen
inputs and recalculate. This native service is offline; it does not
implement the browser's online-source policy. The companion GUI entry is
\texttt{Individual\_tools/14\_Cyclotron\_comparison.command} and its
standalone page is \texttt{cyclotron\_compare}. Module and launcher
names identify the installed application; they do not link to private
source downloads.

The GUI can also be opened with
\texttt{./Run\_private.command\ traplab.py\ gui\ -\/-standalone\ -\/-page\ cyclotron\_compare}.
It offers an English or Chinese switch, one to ten rows, a
frequency-by-row plot, expandable source details, and Auto, TrapLab,
Cloudberry or MRS legacy notation. Native examples include
\texttt{133Cs}, \texttt{133Xe{[}m1{]}} and \texttt{12C:16O2}; full
element-name aliases are not advertised by this native page. Its Save
action creates a fresh \texttt{outputs/cyclotron\_compare/run-N/}
bundle, and Open rows JSON validates all settings before replacing them
and recalculating. The dedicated bundle contains CSV and JSON; figure
saving uses the existing Plot style and Export plots facilities, not an
assumed all-format bundle.

\subsection{Recovery and evidence}\label{recovery-and-evidence}

Resolve ambiguous species by naming isotopes, select an available
nuclear state, correct charge conflicts, and recalculate. If the service
fails, choose saved data explicitly or retry; inputs are retained. For
manual input, fill a finite positive mass and choose its convention.
Empty uncertainty must not be replaced with zero merely to obtain an
error bar.

The web source has 43 focused parser, calculation and UI tests. Eleven
native core and actual BasePage tests passed, including signed charges,
state errors, repeated-isotope uncertainty, units, unknown and zero
errors, settings reopening and invalid-input preservation. Offscreen
construction does not establish native Cocoa export, browser download,
hardware or precision mass-measurement acceptance. The mathematical
result is conditional on the selected mass source and omitted binding
corrections.

\subsection{Primary references}\label{primary-references}

\begin{itemize}
\tightlist
\item
  M. Wang, W. J. Huang, F. G. Kondev, G. Audi and S. Naimi, The AME2020
  atomic mass evaluation II, Chinese Physics C 45, 030003 (2021);
  \href{https://www-nds.iaea.org/amdc/ame2020/mass_1.mas20.txt}{official
  unrounded mass table}. Table flags distinguish estimated values and
  missing quantities.
\item
  F. G. Kondev et al., The NUBASE2020 evaluation of nuclear physics
  properties, Chinese Physics C 45, 030001 (2021),
  \href{https://doi.org/10.1088/1674-1137/abddae}{DOI
  10.1088/1674-1137/abddae};
  \href{https://www-nds.iaea.org/amdc/ame2020/nubase_4.mas20.txt}{official
  state table}.
\item
  IAEA Nuclear Data Services,
  \href{https://www-nds.iaea.org/relnsd/vcharthtml/api_v0_guide.html}{LiveChart
  Data Download API guide}, checked 24 October 2026. This guide
  describes API fields; it does not change the selected AME/NUBASE
  evaluation.
\end{itemize}

\end{document}
